\documentclass[main]{subfiles}


\title[Dynamic Approaches]{Dynamic Approaches}
\subtitle{White Box Optimization}
\author[André Biedenkapp]{André Biedenkapp}
\date{}

\titlebackground*{images/AutoML_LearningToLearn2}

%\institute{}
%\date{}

\begin{document}
	
	\maketitle
%-----------------------------------------------------------------------
%----------------------------------------------------------------------
%-----------------------------------------------------------------------
%----------------------------------------------------------------------
\begin{frame}{Black vs. Gray vs. White Box} 

    \begin{columns}
    \column{0.33\textwidth}
    \begin{itemize}
        \item Black box:
        \begin{itemize}
            \item Choose the input
            \item Observe the final outcome
        \end{itemize}
        \hfill\\
        No control over the box\\
        \hfill\\
        \pause
    \end{itemize}
    \column{0.66\textwidth}
    \begin{itemize}
        \item Gray box (e.g., multi-fidelity):
        \begin{itemize}
            \item Choose the input
            \item Observe intermediate outcome
            \item \ldots
            \item Observe final outcome
        \end{itemize}
        \hfill\\
        Some control through, e.g., early termination
    \end{itemize}
    \end{columns}
    \vspace{6ex}
    \pause
    AutoML as a white-box problem:
    \begin{itemize}
        \item Observe and control \emph{all} details of an algorithm run
        \item Goal: Replace hyperparameter heuristics or even complete algorithm components with learned policies
    \end{itemize}

\end{frame}
%-----------------------------------------------------------------------
%----------------------------------------------------------------------
\begin{frame}{Iterative Optimization Heuristics}
\begin{colorblock}[black]{sintefgrey}{IOHs}
Iterative Optimization Heuristics (IOHs) propose a set of solution candidates in each iteration based on previous evaluations.
\end{colorblock}
\pause
Important Observations:
\begin{itemize}
    \item Many ML and AI algorithms are iterative in nature\pause
    \item A solution candidate refers to a (partially) trained model\\ (e.g., weights of a neural network)\pause
    \item The goal is to improve the quality of successive solution candidates\pause
    \item The main component is a proposal mechanism/heuristic for new solution candidates
\end{itemize}
\end{frame}
%-----------------------------------------------------------------------
%----------------------------------------------------------------------
\begin{frame}{Adaptation of IOHs}
\begin{colorblock}[black]{sintefgrey}{Dynamic Adaptation of Hyperparameters}
The goal is to dynamically adapt hyperparameters based on some feedback from the algorithm.
\end{colorblock}
Examples:\\
\begin{itemize}
    \item Self-adaptation in Evolutionary Algorithms (see, e.g., Step-Size Adaptation in Section 4 of \liturl{Hansen, 2006}{http://www.cmap.polytechnique.fr/~nikolaus.hansen/hansenedacomparing.pdf})
    \item Reduction of the learning rate when the loss stops improving (e.g., \liturl{Pytorchs ReduceLROnPlateau}{https://docs.pytorch.org/docs/stable/generated/torch.optim.lr_scheduler.ReduceLROnPlateau.html\#torch.optim.lr_scheduler.ReduceLROnPlateau})
    \item Adaptive policy update constraint in PPO \liturl{Schulman et al. 2017}{https://arxiv.org/pdf/1707.06347}
\end{itemize}
\end{frame}
%-----------------------------------------------------------------------
%----------------------------------------------------------------------
\begin{frame}{Configuration of IOHs}
\begin{colorblock}[black]{sintefgrey}{Dynamic Algorithm Configuration (DAC)}
The goal of DAC is to learn a \alert{policy} from data that adapts the hyperparameter settings of an IOH.
\end{colorblock}
Examples:\\
\begin{itemize}
    \item Learned Learning Rate adaptation for SGD \liturl{Daniel et al. 2016}{https://ojs.aaai.org/index.php/AAAI/article/view/10187}
    \item Learned mutation operator selection for differential evolution \liturl{Sharma et al. 2019}{https://dl.acm.org/doi/10.1145/3321707.3321813}
    \item Learned interaction frequency for RL \liturl{Biedenkapp et al. 2021}{https://arxiv.org/abs/2106.05262}
\end{itemize}
\end{frame}
%-----------------------------------------------------------------------
%----------------------------------------------------------------------
% \begin{frame}{Tuning of IOHs}
% \begin{colorblock}[black]{sintefgrey}{Learning to Learn (L2L)}
% The goal of L2L is to learn a proposal mechanism/heuristic directly from data. (Not covered in this lecture.)
% \end{colorblock}
% Examples:\\
% \begin{itemize}
%     \item Learned gradient-based optimization \liturl{Andrychowicz et al. 2016}{https://arxiv.org/abs/1606.04474}
%     \item Learned Step-Size adaptation Heuristic for CMA-ES \liturl{Shala et al. 2020}{https://andrebiedenkapp.github.io/assets/pdf/paper/20-PPSN-LTO-CMA.pdf}
%     \item Learning to Meta-Reinforcement Learn via Transformers \liturl{Shala et al. 2025}{https://openreview.net/forum?id=UENQuayzr1}
% \end{itemize}
% \end{frame}
%-----------------------------------------------------------------------
%----------------------------------------------------------------------
\begin{frame}{Main Takeaways} 

%Most important insights from this video:

\begin{itemize}
    \item \textbf{White-box Optimization} The goal is to replace hyperparameter heuristics or even complete algorithm components
with learned policies.
    \item \textbf{IOHs} propose a set of solution candidates in each
iteration based on previous evaluations.
    \item \textbf{Tuning of IOHs} includes adaptation and configuration.
\end{itemize}

 \end{frame}


	

%----------------------------------------------------------------------
%----------------------------------------------------------------------
% \begin{frame}[c]{Planning Slide}
% Likely very close to original slide deck
% \begin{enumerate}
%     \item BB vs GB vs WB
%     \item IOH
% \end{enumerate}
% \end{frame}
% %-----------------------------------------------------------------------

\end{document}
